\subsection{Dock motion in ROV operations}
\label{sec:lit-dock-motion}

During operational practice, the docking target is rarely still in place. For instance, Trslic's docking system has a work-class ROV recovered into a TMS cage suspended from its support vessel, and the cage inherits the vessel's wave-driven motion through the lift wire. Manual TMS docking requires the pilot's full concentration and skill, flying into a cage whose heave, yaw and pendulum swing follow the sea state, while decisions are taken in fractions of a second \cite{trslic2020vision}. The first autonomous TMS docking trials were run in dynamic North Atlantic wave conditions and were motivated by exactly this difficulty, with the objective to widen the operational weather window that the task currently dictates \cite{trslic2020vision}. In some concepts, the dock is not even a fixed structure. A docking-station ROV captures and recharges an AUV in the water column, so both the dock and the AUV are free bodies in the system \cite{morinaga2021development}. Even purpose-built resident docking stations are not still, since they are typically moored or suspended structures rather than rigid foundations. Dynamic analysis of a seabed-moored suspended dock shows it responds measurably to ocean-current forcing transmitted through its mooring line \cite{guo2024dynamic}.

This phenomenon of forcing on underwater object's motion is ordinary ocean physics. Ocean surface waves carry their energy at period of roughly 5 to 20~s (crest-to-crest times). Waves raised by the local wind sit at the short end of this band (5 to 8~s period), while swell arriving from distant storms has longer periods (above ~8~s) \cite{holthuijsen2007waves}. Below the surface, the wave orbital motion decays exponentially with depth at a rate set by the wavelength, while wavelength grows with the square of the period, so short wind sea is filtered out within tens of metres of depth while long-period swell induces wave motion far deeper \cite{holthuijsen2007waves}. A moored dock, therefore, oscillates with a period-dependent amplitude. In a typical sea, the motion surviving at continental-shelf depths is still of order 0.1~m, and because decay depends on period, a dock's depth determines which wave periods reach it. Off the NSW coast, the operating context of this work, the long-term mean sea state is a wave height of 1.6~m at a mean period of 8~s \cite{short1992wave,kinsela2024nearshore}, which places realistic dock motion within the wave band swept in Part~A. The quantitative mapping from period and depth to dock displacement is developed with the experimental design in Section~\ref{sec:method-sim}.

Against this operational picture, the docking control literature almost always models the dock as a fixed pose in the world frame, and where environmental disturbance appears, it is a steady current applied to the vehicle. The exceptions are partial. Trslic et al.\ docked to a heaving TMS cage by servoing on the most recent measured pose, reacting to the motion without having to model it \cite{trslic2020vision}. The same group later showed that TMS heave can be predicted (about 2.5~s ahead with a neuro-fuzzy model), but the prediction came from a pressure sensor mounted on the TMS cage itself, which only covers heave and was evaluated offline instead of in closed-loop docking \cite{trslic2020neuro}. Neither line of work estimates the dock's motion through the vehicle's own camera while both bodies are moving. The estimation problem that configuration raises is taken up in Section~\ref{sec:lit-estimation}.
